% Generated draft: requires mathematical and source review.
\documentclass[11pt]{article}
\usepackage[margin=25mm]{geometry}
\usepackage{fontspec}
\setmainfont{Latin Modern Roman}
\usepackage{amsmath,amssymb,mathtools,unicode-math}
\setmathfont{Latin Modern Math}
\usepackage{enumitem,needspace,xurl,hyperref}
\setlength{\parindent}{0pt}
\setlength{\parskip}{5pt}
\setlength{\emergencystretch}{3em}
\newcommand{\N}{\mathbb{N}}
\newcommand{\Z}{\mathbb{Z}}
\newcommand{\Q}{\mathbb{Q}}
\newcommand{\R}{\mathbb{R}}
\newcommand{\C}{\mathbb{C}}
\newcommand{\F}{\mathbb{F}}
\newcommand{\A}{\mathbb{A}}
\newcommand{\PP}{\mathbb{P}}
\newcommand{\E}{\mathbb{E}}
\newcommand{\eps}{\varepsilon}
\newcommand{\dd}{\,\mathrm{d}}
\begin{document}
\section*{UnsolvedMath ID: 6800027}
\textbf{Source problem code:} catalogue entry
\subsection{Definitions and mathematical statement}
A smooth strongly convex Finsler metric \(F:TM\to[0,\infty)\) is positive away from the zero section, positively homogeneous of degree one in each tangent fiber, and has positive definite fiber Hessian of \(F^2/2\) away from zero. Symmetry is understood in the usual reversible sense,
\[
 F(v)=F(-v).
\]
It is not an assumption that the manifold is a symmetric space. Positive flag curvature means \(K(v,\sigma)>0\) for each tangent 2-plane \(\sigma\) and nonzero flagpole \(v\in\sigma\). A submanifold is totally geodesic when its induced-metric geodesics are also ambient Finsler geodesics.

Ziller asks whether Frankel's theorem extends to symmetric Finsler metrics. In its standard compact totally geodesic intersection form, the proposed extension is: if \((M^n,F)\) is a closed connected reversible Finsler manifold of positive flag curvature, and \(N_1,N_2\subset M\) are nonempty closed totally geodesic submanifolds with
\[
 \dim N_1+\dim N_2\geq n,
\]
then \(N_1\cap N_2\ne\varnothing\).

This supplies the conventional meaning suggested by the source title, ``Totally geodesic submanifolds and positive curvature'', and the supplied summary. The extracted question itself does not state the version of Frankel's theorem, compactness hypotheses, curvature notion, or regularity/embedding conventions. These details must be verified rather than treated as explicit source wording. In particular no minimal-hypersurface version under positive Ricci curvature is substituted.
\subsection{Short English statement}
Does the usual intersection theorem for closed totally geodesic submanifolds in positive curvature remain true for reversible Finsler metrics with positive flag curvature? The exact version of Frankel's theorem intended by the short source question needs confirmation.
\subsection{Sources}
\begin{itemize}
\item Local UnsolvedMath registry, permanent ID 6800027, source code catalogue entry; attributed to Wolfgang Ziller, University of Pennsylvania; proposed year 2018; update 2026-08-06.
\item Morgan and Pansu, \emph{A List of Open Problems in Differential Geometry} (2018), Question environment 27, ``Totally geodesic submanifolds and positive curvature'': \url{https://www.imo.universite-paris-saclay.fr/~pansu/problems_MTDG.tex}. Registry access date 2026-07-29; source-TeX extraction; public source on a coauthor's publication page; present status and redistribution terms need review.
\item The supplied partial/uncertain summary mentions classical Riemannian Frankel theory and unspecified Finsler intersection literature but supplies no exact theorem citation or verified resolution.
\end{itemize}
\end{document}
